\writeSectionFrame{\presentationTitle}{Specification Pattern}
\begin{frame}{Specification Pattern}
	\begin{block}{Definition}
		Specifications machen Geschäftsregeln explizit, 
		wiederverwendbar und kombinierbar.
	\end{block}
\end{frame}


\begin{frame}{Specification Pattern}
	\begin{block}{Grundidee}
		Geschäftsregeln werden in eigene Objekte ausgelagert.
	\end{block}
	\pause
	\begin{itemize}
	  \item Eine Regel = eine Specification
	  \item Regeln sind benannt und wiederverwendbar
	  \item Regeln können kombiniert werden
	\end{itemize}
\end{frame}

\begin{frame}[fragile]{Specification – Grundstruktur}
	\begin{exampleblock}{simple.Specification.java}
		\begin{lstlisting}
public interface Specification<T> {
    boolean isSatisfiedBy(T candidate);
}
		\end{lstlisting}
	\end{exampleblock}
\end{frame}

\begin{frame}[fragile]{Konkrete Specifications}
	\begin{exampleblock}{simple.HasEnoughGold.java}
		\begin{lstlisting}
public class HasEnoughGold
        implements Specification<Player> {

    @Override
    public boolean isSatisfiedBy(Player player) {
        return player.getGold() >= 100;
    }
}
		\end{lstlisting}
	\end{exampleblock}
\end{frame}

\begin{frame}[fragile]{Konkrete Specifications}
	\begin{exampleblock}{simple.IsHighLevel.java}
		\begin{lstlisting}
public class IsHighLevel
        implements Specification<Player> {

    @Override
    public boolean isSatisfiedBy(Player player) {
        return player.getLevel() >= 10;
    }
}
		\end{lstlisting}
	\end{exampleblock}
\end{frame}

\begin{frame}{Kombinieren von Regeln}
	\begin{itemize}
	  \item Regeln lassen sich logisch verknüpfen
	  \item AND, OR, NOT
	  \item Komplexe Logik aus einfachen Bausteinen
	\end{itemize}
\end{frame}

\begin{frame}[fragile]{Composed Specification}
	\begin{exampleblock}{simple.AndSpecification.java}
		\begin{lstlisting}
public class AndSpecification<T>
        implements Specification<T> {
    private final Specification<T> first;
    private final Specification<T> second;

    public AndSpecification(
            Specification<T> first,
            Specification<T> second) {
        this.first = first;
        this.second = second;
    }

    @Override
    public boolean isSatisfiedBy(T candidate) {
        return first.isSatisfiedBy(candidate)
            && second.isSatisfiedBy(candidate);
    }
}
		\end{lstlisting}
	\end{exampleblock}
\end{frame}

\begin{frame}[fragile]{Verwendung}
	\begin{exampleblock}{simple.Main.java}
		\begin{lstlisting}
Specification<Player> canBuyItem =
    new AndSpecification<>(
        new HasEnoughGold(),
        new IsHighLevel()
    );

if (canBuyItem.isSatisfiedBy(player)) {
    // Spieler darf kaufen
}
		\end{lstlisting}
	\end{exampleblock}
\end{frame}

\begin{frame}{Erweiterung der Lösung}
	\begin{itemize}
	  \item verknüpfungs-Spezifikationen als FluentAPI
	  \item Neue abstrakte Basisklasse AbstractSpecification
	  \item AbstractSpecification erzeugt Verknüpfungen and, or, not
	  \item Neue Verknüpfungsklassen OrSpecification, AndSpecification, NotSpecification
	\end{itemize}
\end{frame}

\begin{frame}[fragile]{Oder-Verknüpfung}
	\begin{exampleblock}{fluent.OrSpecification.java}
		\begin{lstlisting}
public class OrSpecification<T> 
		extends AbstractSpecification<T> {
    private final Specification<T> first;
    private final Specification<T> second;

    public OrSpecification(Specification<T> first,
    		Specification<T> second) {
        this.first = first;
        this.second = second;
    }

    @Override
    public boolean isSatisfiedBy(T candidate) {
        return first.isSatisfiedBy(candidate) || 
        	second.isSatisfiedBy(candidate);
    }
}
		\end{lstlisting}
	\end{exampleblock}
\end{frame}

\begin{frame}[fragile]{Abstrakte Basisklasse}
	\begin{exampleblock}{fluent.AbstractSpecification.java}
		\begin{lstlisting}
public abstract class AbstractSpecification<T> 
		implements Specification<T> {
    @Override
    public abstract boolean isSatisfiedBy(T candidate);

    public AbstractSpecification<T> and(Specification<T> other) {
        return new AndSpecification<>(this, other);
    }

    public AbstractSpecification<T> or(Specification<T> other) {
        return new OrSpecification<>(this, other);
    }

    public AbstractSpecification<T> not() {
        return new NotSpecification<>(this);
    }
}
		\end{lstlisting}
	\end{exampleblock}
\end{frame}

\begin{frame}[fragile]{Konkrete Spezifikation}
	\begin{exampleblock}{fluent.HasEnoughGold.java}
		\begin{lstlisting}
public class HasEnoughGold extends 
			AbstractSpecification<Player> {

    @Override
    public boolean isSatisfiedBy(Player player) {
        return player.getGold() >= 100;
    }
}
		\end{lstlisting}
	\end{exampleblock}
\end{frame}

\begin{frame}[fragile]{Konkrete Spezifikation}
	\begin{exampleblock}{fluent.Main.java}
		\begin{lstlisting}
Specification<Player> canBuyItem =
    new HasEnoughGold()
        .and(new IsHighLevel())
        .and(new IsNotCursed());

if (canBuyItem.isSatisfiedBy(player)) {
    System.out.println(player.getName() + " darf kaufen");
} else {
	System.out.println(player.getName() + " darf nicht kaufen");
}
		\end{lstlisting}
	\end{exampleblock}
\end{frame}
